\section{参考文献}
\begingroup
\renewcommand{\refname}{参考文献}
\zihao{-4}\songti
\setlength{\itemsep}{3pt}
\begin{thebibliography}{99}

\bibitem{gupta2008}
GUPTA D, DENTON B. Appointment scheduling in health care: Challenges and opportunities[J]. IIE Transactions, 2008, 40(9): 800--819. DOI: \url{https://doi.org/10.1080/07408170802165880}.

\bibitem{ahmadi2017}
AHMADI-JAVID A, JALALI Z, KLASSEN K J. Outpatient appointment systems in healthcare: A review of optimization studies[J]. European Journal of Operational Research, 2017, 258(1): 3--34. DOI: \url{https://doi.org/10.1016/j.ejor.2016.06.064}.

\bibitem{nhc2020}
国家卫生健康委员会办公厅. 关于进一步完善预约诊疗制度加强智慧医院建设的通知：国卫办医函〔2020〕405号[EB/OL]. (2020-05-21)[2026-08-24]. \url{https://www.nhc.gov.cn/yzygj/c100068/202005/43b2d23ff48448ffae96700bc6eaccd7.shtml}.

\bibitem{chenjuan2023}
陈娟, 李志会, 赵浩, 等. 分时段分项目组多途径预约模式在妇幼专科医院超声检查中的应用[J]. 中国妇幼卫生杂志, 2023, 14(6): 77--80. DOI: \url{https://doi.org/10.19757/j.cnki.issn1674-7763.2023.06.014}.

\bibitem{marynissen2019}
MARYNISSEN J, DEMEULEMEESTER E. Literature review on multi-appointment scheduling problems in hospitals[J]. European Journal of Operational Research, 2019, 272(2): 407--419. DOI: \url{https://doi.org/10.1016/j.ejor.2018.03.001}.

\bibitem{apergi2020}
APERGI L A, BARAS J S, GOLDEN B L, et al. An optimization model for multi-appointment scheduling in an outpatient cardiology setting[J]. Operations Research for Health Care, 2020, 26: 100267. DOI: \url{https://doi.org/10.1016/j.orhc.2020.100267}.

\bibitem{wang2019}
王珊珊, 李金林, 彭春, 等. 不确定服务时间下分布式鲁棒门诊预约调度和排程[J]. 系统工程学报, 2019, 34(4): 566--576.

\bibitem{hrhospital2025}
华容区人民医院. 超声科流程及注意事项[EB/OL]. (2025-08-20)[2026-08-24]. \url{https://www.hbhr.gov.cn/hrqxxgk/xxgkml/ggqsydwxxgk/wsjk/hrqrmyy/202508/t20250821_720866.html}.

\bibitem{zhu2015}
朱顺痣, 王大寒, 何亚男, 等. 基于时间序列模型的医院门诊量分析与预测[J]. 中国科学技术大学学报, 2015, 45(10): 795--803. DOI: \url{https://doi.org/10.3969/j.issn.0253-2778.2015.10.001}.

\bibitem{xiang2009}
向前, 陈平雁. 预测医院门诊量的ARIMA模型构建及应用[J]. 南方医科大学学报, 2009, 29(5): 1076--1078. PMID: 19460744.

\bibitem{luo2017}
LUO L, LUO L, ZHANG X, et al. Hospital daily outpatient visits forecasting using a combinatorial model based on ARIMA and SES models[J]. BMC Health Services Research, 2017, 17: 469. DOI: \url{https://doi.org/10.1186/s12913-017-2407-9}.

\bibitem{lai2025}
LAI C H, LU Y J, CHEN P S. Using data-driven techniques to predict outpatient ultrasound examination time for the multi-clinic outpatient appointment scheduling problem[J]. Communications in Statistics-Simulation and Computation, 2025, 54(9): 3358--3376. DOI: \url{https://doi.org/10.1080/03610918.2024.2349168}.

\bibitem{patrick2008}
PATRICK J, PUTERMAN M L, QUEYRANNE M. Dynamic multipriority patient scheduling for a diagnostic resource[J]. Operations Research, 2008, 56(6): 1507--1525. DOI: \url{https://doi.org/10.1287/opre.1080.0590}.

\bibitem{patrick2007}
PATRICK J, PUTERMAN M L. Improving resource utilization for diagnostic services through flexible inpatient scheduling: A method for improving resource utilization[J]. Journal of the Operational Research Society, 2007, 58(2): 235--245. DOI: \url{https://doi.org/10.1057/palgrave.jors.2602242}.

\bibitem{chen2022}
CHEN P S, CHEN G Y H, LIU L W, et al. Using simulation optimization to solve patient appointment scheduling and examination room assignment problems for patients undergoing ultrasound examination[J]. Healthcare, 2022, 10(1): 164. DOI: \url{https://doi.org/10.3390/healthcare10010164}.

\bibitem{qiao2024}
乔岩, 冉伦, 李金林, 等. 基于两阶段随机规划的远程会诊预约调度问题研究[J]. 中国管理科学, 2024, 32(1): 86--93. DOI: \url{https://doi.org/10.16381/j.cnki.issn1003-207x.2020.1989}.

\bibitem{nhc2022}
国家卫生健康委员会办公厅. 关于印发超声诊断等5个专业医疗质量控制指标（2022年版）的通知：国卫办医函〔2022〕161号[EB/OL]. (2022-05-27)[2026-08-24]. \url{https://www.nhc.gov.cn/yzygj/c100068/202205/164e6556f6d34c33a3344fec9a2c5074.shtml}.

\end{thebibliography}
\endgroup
